% Part: computability
% Chapter: recursive-functions
% Section: bounded-minimization

\documentclass[../../../include/open-logic-section]{subfiles}

\begin{document}

\olfileid{cmp}{rec}{bmi}
\olsection{పరిమిత కనిష్ఠీకరణ}

\begin{explain}
ఏదైనా ధర్మం లేదా సంబంధం~$P$ను నెరవేర్చే అతి చిన్న సంఖ్యగా ఒక
ప్రమేయాన్ని నిర్వచించడం తరచూ ఉపయోగపడుతుంది. $P$ నిర్ణయించదగినదైతే,
$P$ను నెరవేర్చే అతి చిన్న సంఖ్య దొరికేవరకు సాధ్యమైన సంఖ్యలన్నిటినీ
$0$, $1$, $2$, \dots అంటూ ప్రయత్నించడం ద్వారా ఈ ప్రమేయాన్ని
గణించవచ్చు. ఇటువంటి అపరిమిత అన్వేషణ మనలను ఆదిమ పునరావృత్త ప్రమేయాల
పరిధి వెలుపలికి తీసుకెళ్తుంది. అయితే, వేరుగా ఇచ్చిన ఒక అవధికన్నా
\emph{చిన్నదైన} అతి చిన్న సంఖ్యపైనే ఆసక్తి ఉంటే, ఆదిమ పునరావృత్త
పరిధిలోనే ఉంటాం. మరొక విధంగా, కొంచెం సాధారణంగా, మనకు ఆదిమ పునరావృత్త
సంబంధం~$R(x,z)$ ఉందనుకుందాం. $R(x,z)$ నెరవేరే అతి చిన్న $z<y$కు
$x$,~$y$లను పంపే ప్రమేయాన్ని పరిశీలించండి. $R(x,0)$, $R(x,1)$,
\dots, $R(x,y-1)$లను పరీక్షించి దానినీ గణించవచ్చు. అయితే అది ఆదిమ
పునరావృత్త ప్రమేయం ఎందుకు అవుతుంది?
\end{explain}

\begin{prop}
$R(\vec x, z)$ ఆదిమ పునరావృత్త సంబంధమైతే, $R(\vec x, z)$ నెరవేరే
$y$కన్నా చిన్న అతి చిన్న~$z$ ఉంటే దానినీ, లేకపోతే $y$నూ తిరిగి ఇచ్చే
ప్రమేయం $m_R(\vec{x}, y)$ కూడా ఆదిమ పునరావృత్త ప్రమేయమే. ఈ
ప్రమేయం~$m_R$ను
\[
\bmin{z < y}{R(\vec{x}, z)},
\]
అని రాస్తాం.
\end{prop}

\begin{proof}
$R(\vec{x}, z)$ నెరవేరే $z < 0$ ఉండదు; ఎందుకంటే $z < 0$ అసలే
లేదు. అందువల్ల $m_R(\vec x, 0) = 0$.

అవధి $y+1$ రూపంలో ఉన్నప్పుడు మూడు సందర్భాలు ఉన్నాయి:
\begin{enumerate}
\item $R(\vec{x}, z)$ నెరవేరే $z < y$ ఉంది; అప్పుడు
$m_R(\vec{x}, y+1) = m_R(\vec{x}, y)$.
\item అటువంటి~$z<y$ లేదు, కానీ $R(\vec{x}, y)$ నెరవేరుతుంది; అప్పుడు
$m_R(\vec{x}, y+1) = y$.
\item $R(\vec{x}, z)$ నెరవేరే $z < y+1$ ఏదీ లేదు; అప్పుడు
$m_R(\vec{x}, y+1) = y+1$.
\end{enumerate}
\sourcecorrection{OLTECRREM-002}{మూడవ సందర్భంలో ఆంగ్ల మూలం మొదటి ఆర్గుమెంటును $\vec x$ నుంచి సంబంధానికి బంధితమైన $\vec z$గా మార్చింది; మిగతా నిర్వచనంతో సరిపడే $m_R(\vec x,y+1)$ను పునరుద్ధరించాం.}
కాబట్టి $m_R(\vec x, 0)$ను ఆదిమ పునరావృత్తితో ఇలా నిర్వచించవచ్చు:
\begin{align*}
m_R(\vec{x}, 0) & = 0\\
m_R(\vec{x}, y+1) & =
\begin{cases}
m_R(\vec{x}, y) & \text{$m_R(\vec x,y)\neq y$ అయితే}\\
y & \text{$m_R(\vec x,y)=y$ మరియు $R(\vec x,y)$ అయితే}\\
y+1 & \text{ఇతర సందర్భాల్లో.}
\end{cases}
\end{align*}
$R(\vec x, z)$ నెరవేరే $z<y$ ఉండటం, $m_R(\vec x, y) \neq y$ కావడం
సమానార్థకాలని గమనించండి.
\end{proof}

\begin{prob}
$R(\vec x, z)$ ఆదిమ పునరావృత్త సంబంధం అనుకుందాం. $R(\vec{x}, z)$
నెరవేరే $y$కన్నా చిన్న అతి చిన్న~$z$ ఉంటే దానినీ, లేకపోతే $0$నూ
తిరిగి ఇచ్చే ప్రమేయం $m'_R(\vec{x}, y)$ను $\Char{R}$ నుంచి ఆదిమ
పునరావృత్తితో నిర్వచించండి.
\end{prob}

\end{document}
